\documentclass{beamer}

\usepackage[utf8]{inputenc}

\input{../helper}

\title{Introdução à Computação \\ recursão}
\author{Alexandre Rademaker}
\date{}

\begin{document}

\begin{frame}
 \maketitle
\end{frame}

\begin{frame}[fragile,allowframebreaks]{Recursão}

  Imagine queremos imprimir blocos como

  \vspace{1cm}

  \begin{lstlisting}[basicstyle=\ttfamily]
    #
    ##
    ###
    ####
  \end{lstlisting}

  \framebreak

  Nosso código para isso pode ser

  \lstinputlisting[style=tinyc,caption=src/draw0.c]{src/draw0.c}

  \begin{lstlisting}[basicstyle=\ttfamily\scriptsize]
    > make draw0 && ./draw0 4
    #
    ##
    ###
    ####
  \end{lstlisting}

  \framebreak

  Mas observe, pirâmides são estruturas recursivas!

  uma pirâmide de altura 4 é na verdade uma pirâmide de altura 3 com
  uma fileira extra de 4 blocos adicionados.

  e uma pirâmide de altura 3 é uma pirâmide de altura 2 com uma
  fileira extra de 3 blocos.

  e uma pirâmide de altura 2 é uma pirâmide de altura 1 com uma
  fileira extra de 2 blocos.

  e, finalmente, uma pirâmide de altura 1 é uma pirâmide de altura 0
  (sem blocos) com uma linha de 1 bloco adicionada.

  \framebreak

  \lstinputlisting[style=normalc,caption=src/draw1.c]{src/draw1.c}

\end{frame}


\begin{frame}[fragile,allowframebreaks]{Recursão eficiente}

  \begin{lstlisting}[style=ssc,caption=src/math.c]
    long int fib (int n) {
      if( n == 1 || n == 2)
        return 1;
      else
        return fib(n - 1) + fib(n - 2);
    }
  \end{lstlisting}

  Qual o problema com esta solução para Fibonacci:
  \[ F_n = F_{n - 1} + F_{n - 2} \] onde $F_1 = F_2 = 1$?

  \framebreak

  \includegraphics[width=.9\textwidth]{fibo-rec.png}

  Repetição desnecessária de computações. Como corrigir?

  \framebreak

  \begin{align*}
    fib_{aux}(4, 1, 1) & = fib_{aux}(3, 1, 2) \\
    fib_{aux}(3, 1, 2) & = fib_{aux}(2, 2, 3) \\
    fib_{aux}(2, 2, 3) & = fib_{aux}(1, 3, 5) \\
    fib_{aux}(1, 3, 5) & = 3 \\
  \end{align*}

  \begin{lstlisting}[style=ssc,caption=src/math.c]
    long int fib_aux(int n, int a, int b) {
      return (--n > 0) ? fib_aux(n, b, a + b) : a ;
    }

    long int fib(int n) {
      return fib_aux(n, 1, 1);
    }
  \end{lstlisting}
  
\end{frame}


\begin{frame}[fragile,allowframebreaks]{Busca Binária}

  Podemos agora repensar na função de busca binária!

  Considerando um array ordenado, como \emph{recursivamente} verificar
  se um elemento N está no array?

  A cada passo, queremos buscar em um novo array, alguma metade do
  array inicial.

  \framebreak

  \begin{lstlisting}[style=ssc,caption=src/numbers3.c]
    int bin_search_aux(int a[], int x, int i, int j) {
      if (j < i) return -1;

      int k = (j + i) / 2;
      if (a[k] == x)
       return k;
      else if (a[k] < x)
       return bin_search_aux(a, x, k + 1, j);
      else
       return bin_search_aux(a, x, i, k - 1);
    }

    int bin_search(int a[], int x, int n) {
      return bin_search_aux(a, x, 0, n - 1);
    }
  \end{lstlisting}

  Como \textbf{debugar}? Com alguns \textbf{printf}\ldots

  \framebreak

  De uma galáxia muito muito distante\ldots
  \vspace{.5cm}

  \begin{lstlisting}[language=Lisp,style=code,basicstyle=\ttfamily\tiny]
 (defun binary-search (value array &optional (low 0) (high (1- (length array))))
     (if (< high low)
         nil
         (let ((middle (floor (+ low high) 2)))
           (cond ((> (aref array middle) value)
                  (binary-search value array low (1- middle)))
                 ((< (aref array middle) value)
                  (binary-search value array (1+ middle) high))
                 (t middle)))))
   \end{lstlisting}

   \begin{lstlisting}[language=lisp,style=code,basicstyle=\ttfamily\tiny]
   CL-USER> (trace binary-search)
     (BINARY-SEARCH)
     
   CL-USER> (binary-search 4 #(-31 0 1 2 2 4 65 83 99 782))
    0: (BINARY-SEARCH 4 #(-31 0 1 2 2 4 65 83 99 782))
     1: (BINARY-SEARCH 4 #(-31 0 1 2 2 4 65 83 99 782) 5 9)
      2: (BINARY-SEARCH 4 #(-31 0 1 2 2 4 65 83 99 782) 5 6)
      2: BINARY-SEARCH returned 5
     1: BINARY-SEARCH returned 5
    0: BINARY-SEARCH returned 5
    5
  \end{lstlisting}

  \framebreak

  \begin{center}
    \includegraphics[width=.9\textwidth]{lisp_cycles.png}
  \end{center}
  Fonte \url{https://xkcd.com/297/}

\end{frame}


\end{document}

%%% Local Variables:
%%% mode: latex
%%% TeX-master: t
%%% End:
